The goal of this exercise is to write a \texttt{Library} class that represents a library holding several \texttt{Book} objects. It uses the \texttt{Book} class of Exercise 3, which is provided.

\subsection{Implementation}

Write a \texttt{Library} class with the following requirements:

\begin{itemize}

    \item \textbf{Attributes}:

        \begin{itemize}

            \item \texttt{String name}: the name of the library.

            \item \texttt{Book[] books}: the books of the library. The array is filled from its first position, so the empty positions, all at the end, hold \texttt{null}.

            \item \texttt{int currentBookIndex}: the position of the book currently selected. It starts at -1, which means that no book is selected.

        \end{itemize}

    \item \textbf{Methods}:

        \begin{itemize}

            \item \textbf{Constructor}: receives the name of the library and the number of books it has room for at first. The library starts without books.

            \item \texttt{addBook()}: adds a book in the first empty position of the array. If the array is full, it first has to be replaced by one twice as long: since an array cannot change its size, create a new, longer array and copy the books into it, in the same positions.

            \item \texttt{getCurrentBook()}: returns the book currently selected, or \texttt{null} if none is.

            \item \texttt{nextBook()}: selects the next book and returns it. If there are no more books, it returns \texttt{null} and the selection does not change.

            \item \texttt{previousBook()}: selects the previous book and returns it. If there is no book before the current one, or no book is selected, it returns \texttt{null} and the selection does not change.

            \item \texttt{describe()}: returns a \texttt{String} with the data of the library, including its name and its number of books, but not the data of the books themselves.

        \end{itemize}

\end{itemize}

\subsection{Recommendations}

\begin{itemize}

    \item \textbf{Incremental implementation}: write and test each method on its own before moving on to the next one, making sure each part works before building the rest of the class.

\end{itemize}

\subsection{Evaluation and testing}

To complete the exercise, write a \texttt{Main} class with a \texttt{main} method that performs the following tests:

\begin{enumerate}

    \item \textbf{Class functionality}:

        \begin{itemize}

            \item Write the \texttt{Library} class with all the attributes, constructors and methods described.

        \end{itemize}

    \item \textbf{Testing in the} \texttt{main} \textbf{method}:

        \begin{itemize}

            \item Create a library and several books.

            \item Add the books to the library, with a starting size small enough for the array to have to double.

            \item Go through the library with the navigation methods, in both directions, and check that reaching either end is handled properly.

        \end{itemize}

    \item \textbf{Clean code and documentation}:

        \begin{itemize}

            \item Make sure the code is clean and legible, with explanatory comments where necessary.

        \end{itemize}

\end{enumerate}
